\chapter{The linearized operator of a single transition}\label{sec:phase-linearized-transition}
We use $q(t)=\tanh(t/\sqrt2)$, $\psi=q'>0$, $V(t)=W''(q(t))=2-3\operatorname{sech}^2(t/\sqrt2)$ and
\[
 \mathscr L=-\frac{d^2}{dt^2}+V(t),\qquad
 Q(f)=\int_{\mathbb R}(|f'|^2+Vf^2),\qquad
 \|\psi\|_2^2=c_W=\frac{2\sqrt2}{3}.
\]
The positive-gap and kernel questions are posed in \cite{ChAC}, Appendix B, Problem 2, pp.~28--29. We compute the sharp constant for this particular quartic potential. The ground-state identity of Proposition \ref{prop:phase-ground-state} is already available; the new point is coercivity after removing translation.

\begin{lemma}[Two first-order factorizations]\label{lem:phase-two-factorizations}
Set
\[
 \mathcal A=\frac d{dt}+\sqrt2\tanh(t/\sqrt2),\qquad
 \mathcal B=\frac d{dt}+\frac1{\sqrt2}\tanh(t/\sqrt2).
\]
Their formal $L^2$ adjoints reverse the derivative sign, and on smooth functions
\[
 \mathscr L=\mathcal A^*\mathcal A,\qquad
 \mathcal A\mathcal A^*=\mathcal B^*\mathcal B+\frac32,
 \qquad \mathcal B\mathcal B^*=-\frac{d^2}{dt^2}+\frac12.
\]
For $f\in H^1(\mathbb R)$, $Q(f)=\|\mathcal A f\|_2^2$. Its zero space is exactly $\operatorname{span}\{\psi\}$.
\end{lemma}
\begin{proof}
For $D=d/dt$ and a real function $b$, $(D+b)^*(D+b)=-D^2+b^2-b'$ and $(D+b)(D+b)^*=-D^2+b^2+b'$. Insert $b=\sqrt2\tanh(t/\sqrt2)$ and then $b=\tanh(t/\sqrt2)/\sqrt2$, using $\tanh^2+\operatorname{sech}^2=1$. Integration by parts proves the quadratic identity for compactly supported smooth functions. Boundedness of the coefficients and density extend it to $H^1$.

Since $\psi'/\psi=-\sqrt2\tanh(t/\sqrt2)$, $\mathcal A\psi=0$. If $Q(f)=0$, the weak first-order equation gives $(f/\psi)'=0$ on every bounded interval, where $\psi$ is smooth and bounded away from zero. The constants agree on overlaps, so $f=c\psi$ globally. Conversely these functions lie in $H^1$ and have zero quadratic form.
\end{proof}

\begin{proposition}[The bounded kernel on the line]\label{prop:phase-bounded-normal-kernel}
Every bounded $C^2$ solution of $\mathscr Lf=0$ on $\mathbb R$ is a constant multiple of $\psi$. More generally, all solutions have the form
\[
 f(t)=\psi(t)\left(c_1+c_2\int_0^t\psi(s)^{-2}ds\right).
\]
The second independent solution grows exponentially as $t\to+\infty$ and is unbounded.
\end{proposition}
\begin{proof}
Write $f=\psi v$ and use $\mathscr L\psi=0$. The equation becomes $(\psi^2v')'=0$, yielding the formula. Since $\psi(t)=2\sqrt2 e^{-\sqrt2t}(1+o(1))$ as $t\to+\infty$, its second solution is asymptotic to a nonzero constant times $e^{\sqrt2t}$. A bounded solution therefore has $c_2=0$. Conversely $\psi$ is bounded. This proves the one-dimensional assertion in \cite{ChAC}, Problem 2(b), without assuming that a bounded solution is square integrable in advance.
\end{proof}

\begin{lemma}[An adjoint primitive for compact orthogonal data]\label{lem:phase-adjoint-primitive}
If $f\in C_c^\infty(\mathbb R)$ satisfies $\int f\psi=0$, then
\[
 g(t)=-\frac1{\psi(t)}\int_{-\infty}^t f(s)\psi(s)ds
\]
is smooth and compactly supported, $\mathcal A^*g=f$, and
\[
 Q(f)-\frac32\|f\|_2^2
 =\|\mathcal B^*\mathcal B g\|_2^2+\frac32\|\mathcal B g\|_2^2.
\]
Moreover $\|g\|_2^2\leq(2/3)\|f\|_2^2$.
\end{lemma}
\begin{proof}
The integral vanishes to the left of the support, and orthogonality makes it vanish to the right as well. Differentiate $\psi g$ to obtain $\mathcal A^*g=f$. Put $T=\mathcal A\mathcal A^*=\mathcal B^*\mathcal B+3/2$. Then
$Q(f)=\|Tg\|_2^2$ and $\|f\|_2^2=\langle Tg,g\rangle$.
Expand $T$ in these two expressions and integrate by parts; subtracting $3\langle Tg,g\rangle/2$ gives exactly the asserted sum. The lower bound $\langle Tg,g\rangle\geq(3/2)\|g\|_2^2$ gives the final estimate. All integrations here have compact support, so no endpoint behavior is assumed.
\end{proof}

\begin{theorem}[The exact gap away from translation]\label{thm:phase-sharp-normal-gap}
If $f\in H^1(\mathbb R)$ and $\int f\psi=0$, then
\[
 Q(f)\geq\frac32\|f\|_2^2.
\]
Equality holds precisely for multiples of
$\zeta(t)=\operatorname{sech}(t/\sqrt2)\tanh(t/\sqrt2)$.
Thus $3/2$ is the optimal gap. For arbitrary $f\in H^1$,
\[
 Q(f)\geq\frac32\left(\|f\|_2^2-
 \frac{|\int f\psi|^2}{c_W}\right).
\]
The existence of a positive gap is the assertion in \cite{ChAC}, Problem 2(c); the factorization determines its exact value here.
\end{theorem}
\begin{proof}
First apply Lemma \ref{lem:phase-adjoint-primitive} to compact orthogonal tests. For an arbitrary orthogonal $f$, approximate in $H^1$ by compact smooth functions and subtract a coefficient tending to zero times a fixed compact smooth function having nonzero pairing with $\psi$. This produces orthogonal $f_j\to f$ in $H^1$, proving the inequality by continuity.

Let $g_j$ be their adjoint primitives. Applied to differences, the last estimate of the lemma makes $g_j$ Cauchy in $L^2$. Since $g_j'=\sqrt2\tanh(t/\sqrt2)g_j-f_j$, the derivatives are also Cauchy in $L^2$. Thus $g_j\to g$ in $H^1$ and $\mathcal A^*g=f$. If equality holds for $f$, the sum identity implies $\mathcal B g_j\to0$ in $L^2$, hence $\mathcal B g=0$. Solving this first-order equation gives $g=c\chi$, where $\chi=\operatorname{sech}(t/\sqrt2)$. Direct differentiation gives $\mathcal A^*\chi=(3/\sqrt2)\zeta$, so equality forces the stated form. Conversely $\mathcal B\chi=0$ and the factorization give $\mathscr L\zeta=(3/2)\zeta$. This function is in $H^1$ and is orthogonal to the even function $\psi$, so it realizes equality.

For general $f$, subtract its $L^2$ projection onto $\psi$. The quadratic form is unchanged because $\mathcal A\psi=0$, while the squared norm loses exactly the displayed projection term. This proves the last formula.
\end{proof}

\begin{proposition}[All square-integrable eigenfunctions of the normal operator]\label{prop:phase-normal-point-spectrum}
If a nonzero real $f\in H^2(\mathbb R)$ satisfies $\mathscr Lf=\lambda f$, then either $\lambda=0$ and $f$ is a multiple of $\psi$, or $\lambda=3/2$ and $f$ is a multiple of $\zeta$. This classifies square-integrable eigenfunctions; it is not a statement about the full continuous spectrum.
\end{proposition}
\begin{proof}
The quadratic identity gives $\lambda\geq0$. At zero the kernel was identified in Lemma \ref{lem:phase-two-factorizations}. Suppose $\lambda>0$ and set $g=\mathcal A f$. It is nonzero, and intertwining yields
$(\mathcal A\mathcal A^*)g=\lambda g$. Set $h=\mathcal B g$. The second factorization gives
\[
 -h''+2h=\lambda h.
\]
These operations are justified in $L^2$: the eigen-equation and bounded derivatives of $V$ imply successively that all derivatives of $f$ lie in $L^2$, starting from $f\in H^2$ and differentiating $f''=(V-\lambda)f$.

No nonzero solution of the displayed constant-coefficient equation belongs to $L^2(\mathbb R)$. For $\lambda<2$ its solutions are linear combinations of opposite exponentials, for $\lambda=2$ affine functions, and for $\lambda>2$ sines and cosines. Each nonzero combination fails square integrability at one or both ends. Hence $h=0$, so $g=c\chi$ and the nonzero equation for $g$ forces $\lambda=3/2$. Finally $\lambda f=\mathcal A^*g$ gives the stated eigenfunction. The converse eigen-equations have already been checked.
\end{proof}

\begin{theorem}[A normal inverse after fixing translation]\label{thm:phase-orthogonal-normal-inverse}
For every $F\in L^2(\mathbb R)$ with $\int F\psi=0$, there is a unique $u\in H^2(\mathbb R)$ satisfying
\[
 \mathscr Lu=F,\qquad \int u\psi=0.
\]
With $\|u\|_{H^2}^2=\|u\|_2^2+\|u'\|_2^2+\|u''\|_2^2$, one has
\[
 \|u\|_2\leq\frac23\|F\|_2,\qquad
 \|u\|_{H^2}\leq\sqrt7\|F\|_2.
\]
Orthogonality of $F$ is necessary for any $H^2$ solution. This supplies an elementary normal-direction solvability statement for the linearized problem discussed in \cite{ChAC}, Problems 2 and 3; it is not a uniform inverse on a curved or noncompact interface.
\end{theorem}
\begin{proof}
On the closed subspace $\mathcal H=\{u\in H^1:\int u\psi=0\}$, the gap and $V\geq-1$ give
\[
 \|u\|_2^2\leq\frac23Q(u),\qquad
 \|u'\|_2^2\leq Q(u)+\|u\|_2^2\leq\frac53Q(u).
\]
Thus $Q$ is an inner product with norm equivalent to the complete $H^1$ norm on this subspace. Hilbert representation, Lemma \ref{lem:hilbert-functional-representation}, applied to $v\mapsto\int Fv$, gives a unique $u\in\mathcal H$ with $Q(u,v)=\int Fv$ for every $v\in\mathcal H$. Both sides vanish on $\psi$, so decomposing a general $H^1$ test into its orthogonal part and its multiple of $\psi$ proves the equation weakly against all tests. Since $u''=Vu-F\in L^2$, $u\in H^2$.

Testing with $u$ gives $Q(u)\leq\|F\|_2\|u\|_2$ and hence the first estimate. It also gives $Q(u)\leq(2/3)\|F\|_2^2$, so $\|u'\|_2^2\leq(10/9)\|F\|_2^2$. Since $|V|\leq2$, $\|u''\|_2\leq2\|u\|_2+\|F\|_2\leq(7/3)\|F\|_2$. Adding the three squared estimates gives the constant seven. Uniqueness follows from the gap. Finally any $H^2$ solution paired with $\psi$ has $\int F\psi=Q(u,\psi)=0$, proving necessity.
\end{proof}

\begin{corollary}[The gap at a physical interface scale]\label{cor:phase-scaled-normal-gap}
For $\varepsilon>0$, if $f\in H^1(\mathbb R)$ and $\int f(t)\psi(t/\varepsilon)dt=0$, then
\[
 \int\left(\varepsilon|f'|^2+\varepsilon^{-1}V(t/\varepsilon)f^2\right)dt
 \geq\frac3{2\varepsilon}\int f^2dt.
\]
The operator $-\varepsilon^2d^2/dt^2+V(t/\varepsilon)$ instead has gap $3/2$, while $-d^2/dt^2+\varepsilon^{-2}V(t/\varepsilon)$ has gap $3/(2\varepsilon^2)$. These are different normalizations of the same one-dimensional estimate.
\end{corollary}
\begin{proof}
Use $s=t/\varepsilon$ and apply Theorem \ref{thm:phase-sharp-normal-gap} to $f(\varepsilon s)$. The energy on the left becomes its unscaled quadratic form, while its squared norm is $\varepsilon^{-1}\int f^2$. Multiplying the resulting inequality by $\varepsilon$ or by $\varepsilon^{-1}$ gives the other two normalizations. Equality is attained by the correspondingly rescaled function $\zeta$, so the constants remain sharp.
\end{proof}
